\label{sec:results}

% Replace placeholders with final figures/tables before submission.
% Export plots from notebooks or training logs; keep paths under docs/report/semester-project/figures/

\subsection{Baseline overflight}
Figure~\ref{fig:baseline-overflight} shows the reference baseline pass with seeded clouds along the target corridor (notebook \texttt{07-baseline-overflight.ipynb}).

\begin{figure}[htbp]
  \centering
  % \includegraphics[width=0.85\linewidth]{figures/baseline-overflight-frame.png}
  \fbox{\parbox[c][4cm][c]{0.85\linewidth}{\centering Placeholder: export a representative frame or metric plot.}}
  \caption{Baseline overflight scenario (cloud-aware corridor).}
  \label{fig:baseline-overflight}
\end{figure}

\subsection{Training progression}
Report learning curves (episode return, capture count, evaluation rollouts) for the final model checkpoint.
State the random seed policy and number of training episodes.

\subsection{Policy vs.\ baseline}
Present a compact comparison table.
Example headings: controller, mean episode return, captures per orbit, mean quality, safety terminations.

\begin{table}[htbp]
  \centering
  \caption{Controller comparison (fill with final numbers).}
  \label{tab:controller-comparison}
  \begin{tabular}{lcccc}
    \toprule
    Controller & Mean return & Captures & Mean quality & Terminations \\
    \midrule
    Baseline   & --- & --- & --- & --- \\
    MPO (final)& --- & --- & --- & --- \\
    \bottomrule
  \end{tabular}
\end{table}
